It is clear that (viii) follows from (vii), for $\pi=\pi''\circ p$, where $\pi'':G''\to S$ denotes the structural morphism.

Let us show (ix). By (i), $p$ is an epimorphism; since $G$ is reduced, $p$ factors through the immersion $G''_{\mathrm{red}}\to G''$, which is therefore also an epimorphism, hence an isomorphism (IV 4.4.4).

Let us show (x). If $G''$ is separated over $S$, then $\varepsilon''$ is a closed immersion, by EGA I, 5.4.6. On the other hand, we saw in (ii) that $\varepsilon''$ is a closed immersion if and only if $u$ is one. Finally, if $u$ is a closed immersion, then so is $\delta\circ(\mathrm{id}_G\times u):G\times_S G'\to G\times_S G$; thus, by Lemma 9.2.1 below, $G''$ is separated over $S$.

Assertion (xi) follows from (vii) and EGA IV$_2$, 2.2.13.

Assertion (xii) follows from (vii), EGA IV$_4$, 17.7.5 and 17.7.7, and the fact that, since $p$ is universally open, for every $s\in S$, if the underlying space of $G_s$ is discrete, the same holds for the underlying space of $G''_s$.

Finally, assertion (xiii) follows from (vii), (viii), and EGA IV$_4$, 17.7.5.

\begin{lemma}{9.2.1}\label{VIB.9.2.1}
Let $X$ be an $S$-scheme and $R$ an equivalence relation defined on $X$ by the monomorphism $v:R\to X\times_S X$. Suppose $R$ effective. Then:

(i) $v$ is an immersion, and it is a closed immersion if $Y=X/R$ is separated.

(ii) Suppose moreover that $Y$ represents the (fpqc) sheaf quotient of $X$ by $R$\nde{As pointed out by O. Gabber, this is used in the proof and must be inserted among the hypotheses.} and that $v$ is a closed immersion. Then $Y=X/R$ is separated over $S$.
\end{lemma}
Recall (IV Def. 3.3.2) that the hypothesis ``$R$ effective'' means that there is a morphism of $S$-schemes $p:X\to Y$ such that the natural morphism $R\to X\times_Y X$ is an isomorphism. One deduces (EGA I, 5.3.5) the following cartesian diagram:
\[
\xymatrix{R\ar[r]^v\ar[d]&X\times_S X\ar[d]^{p\times p}\\Y\ar[r]_{\Delta_{Y/S}}&Y\times_S Y.}
\]
Then, since $\Delta_{Y/S}$ is an immersion (EGA I, 5.3.9), the same holds for $v$. Similarly, if $Y$ is separated over $S$, $\Delta_{Y/S}$ is a closed immersion, so the same holds for $v$.

Conversely, suppose $v$ a closed immersion and $Y$ represents the (fpqc) sheaf quotient of $X$ by $R$. Then $p$ is covering for the (fpqc) topology (IV 4.4.3), and therefore $p\times p$ is too (by base change, $p\times\mathrm{id}_X$ and $\mathrm{id}_Y\times p$ are covering, and thus also their composite $p\times p$). Hence, by (fpqc) descent (cf. EGA IV$_2$, 2.7.1), $\Delta_{Y/S}$ is a closed immersion, i.e.\ $Y$ is separated over $S$.

\begin{remark}{9.2.2}\label{VIB.9.2.2}
Under the general hypotheses of 9.2, if one supposes $G'$ flat and locally of finite presentation over $S$, then $p$ is covering for the (fppf) topology,\nde{(fpqc) has been corrected to (fppf).} by 9.2 (vii), so assertions (vii) and (viii) of 9.2 can be extended to properties $\mathscr P$ local for the (fppf) topology.
\end{remark}

\begin{remark}{9.3}\label{VIB.9.3}
a) The question whether a quotient $G/G'$ is representable or not is often delicate; in this seminar we prove representability of certain particular quotients.

In general, to be able to assert that the quotient $G/G'$ is representable, it does not suffice to suppose $G$ and $G'$ of finite presentation over $S$ and $G'$ flat over $S$. Indeed, suppose moreover $G$ smooth with connected fibers. In this case, if $G/G'$ is a scheme, it is separated, by Corollary 5.4, and therefore $G'\hookrightarrow G$ is a closed immersion, by 9.2 (x); consequently, if $G'$ is not closed in $G$, then $G/G'$ is not representable.

To obtain such a counterexample, one may take for $S$ the spectrum of a discrete valuation ring, and put $G=\Gm{}_{,S}$. Consider moreover an integer $n>1$, invertible on $S$; then $\bmu_n=\Ker(G\xrightarrow{n}G)$ is a closed subgroup of $G$ étale over $S$ (cf. VIIA\nde{See VIIA, 8.4 or VIII, 2.1.}). Let $G'$ be the open subgroup of $\bmu_n$ obtained by removing from $\bmu_n$ the closed part of the closed fiber of $\bmu_n$ complementary to the origin. Then $G'$ is not closed in $G$, so $G/G'$ is not representable. (One can also construct such examples where $G'$ is smooth with connected fibers.)

b) On the other hand, it is not excluded that $G/G'$ is representable when $G$ and $G'$ are of finite presentation over $S$, and $G'$ is flat over $S$ and closed in $G$.\begingroup\renewcommand{\thefootnote}{*}\addtocounter{footnote}{-1}\footnote{This is too optimistic, as M. Raynaud shows in his thesis (loc.\ cit.\ X 14).}\endgroup\nde{Remark (*) refers to counterexample X.14 in [Ray70a]. The base there is regular local of dimension 2.} Under these hypotheses, one knows that $G/G'$ is representable in the following special cases:

$1^\circ$ -- $S$ is the spectrum of an artinian ring (cf. VIA 3.2 and 3.3.2).

$2^\circ$ -- $G'$ is proper over $S$ and $G$ quasi-projective over $S$ (cf. V 7.1).

$3^\circ$ -- $S$ is locally noetherian of dimension $1$ (cf. [An73], Th. 4.C).
\end{remark}

\sgasection{10}{Passage to the projective limit in group schemes and schemes with a group of operators}
\numpara{10.0}\label{VIB.10.0}
Recall the essential result of EGA IV$_3$, \S\,8.8. Suppose the following situation given: $S_0$ a quasi-compact and quasi-separated scheme, $I$ an increasing filtered preordered set, $(\mathscr A_i)_{i\in I}$ an inductive system of quasi-coherent commutative $\OO_{S_0}$-algebras, $\mathscr A=\varinjlim\mathscr A_i$, $S_i=\Spec\mathscr A_i$ for $i\in I$, and $S=\Spec\mathscr A$;\nde{Note that $S$, being affine over $S_0$, is therefore quasi-compact and quasi-separated, cf. N.D.E. (92) below.} then the category of $S$-schemes of finite presentation is determined up to equivalence by the categories of $S_i$-schemes of finite presentation, the functors between these categories $\rho_{ji}:X_i\mto X_i\times_{S_i}S_j$ for $i\leq j$, and the transitivity isomorphisms $\rho_{kj}\circ\rho_{ji}\isomto\rho_{ki}$.

Let us be more precise. Given $j\in I$, and an $S_j$-scheme of finite presentation $X_j$, we shall put, for every $i\in I$ such that $i\geq j$, $X_i=X_j\times_{S_j}S_i$, and $X=X_j\times_{S_j}S$. Then (EGA IV$_3$, 8.8.2):

(i) Given $j\in I$, and two $S_j$-schemes of finite presentation $X_j$ and $Y_j$, the canonical map from $\varinjlim_{i\geq j}\Hom_{S_i}(X_i,Y_i)$ to $\Hom_S(X,Y)$ is bijective.

(ii) For every $S$-scheme of finite presentation $X$, there is an index $j\in I$, an $S_j$-scheme of finite presentation $X_j$ and an $S$-isomorphism $X\isomto X_j\times_{S_j}S$.

One concludes (EGA IV$_3$, 8.8.3) that, whenever one has a diagram $D$ involving finitely many objects and arrows of the category of $S$-schemes of finite presentation, one can find an index $i\in I$ and a diagram $D_i$ in the category of $S_i$-schemes of finite presentation, such that the diagram $D$ comes, up to isomorphism, from the diagram $D_i$ by base change $S\to S_i$. One can even find $i$ and $D_i$ such that every cartesian square of $D$ comes from a cartesian square of $D_i$.

\numpara{10.1}\label{VIB.10.1}
Moreover, many common properties of a morphism, stable under base change, possess the following property:
\begin{quote}
Let $u:X\to Y$ be an $S$-morphism between $S$-schemes of finite presentation; by 10.0 it comes by base change from an $S_j$-morphism $u_j:X_j\to Y_j$ between $S_j$-schemes of finite presentation; then, for $u$ to have property $\mathscr P$, it is necessary and sufficient that there be $i\geq j$ such that $u_i=u_j\times_{S_j}S_i$ has property $\mathscr P$.
\end{quote}
This is the case when $\mathscr P$ is one of the following properties of a morphism: being separated, surjective, radicial, affine, quasi-affine, finite, quasi-finite, proper, projective, quasi-projective, an isomorphism, a monomorphism, an immersion, an open immersion, a closed immersion (EGA IV$_3$, 8.10.5), flat (EGA IV$_3$, 11.2.6), smooth, unramified or étale (EGA IV$_4$, 17.7.8).\nde{The word ``flat'' has been added, and 17.7.6 corrected to 17.7.8.}

Note that this still holds when $\mathscr P$ is the property of being covering for the (fppf) topology; indeed, given two $S$-schemes of finite presentation $X$ and $Y$, and an $S$-morphism $u:X\to Y$, it follows from IV, 6.3.1 (i)\nde{And from the fact that $Y$, being of finite presentation over $S$, is quasi-compact.} that, for $u$ to be covering for the (fppf) topology, it is necessary and sufficient that there be an $S$-scheme $Z$ and an $S$-morphism $v:Z\to Y$ faithfully flat and of finite presentation which factors through $u$.

The aim of this section 10 is to give variants of results of this kind for the category of $S$-groups of finite presentation, that of $S$-schemes with a group of operators, and for certain properties of group monomorphisms (being invariant, central with representable quotient sheaf, etc.).

The first two results of this type are the following. (In nos.\ 10.2 to 10.9 below, we retain the notation introduced in 10.0.)

\begin{lemma}{10.2}\label{VIB.10.2}
Let $G_j$ and $H_j$ be two $S_j$-groups of finite presentation; put, for every $i\geq j$, $G_i=G_j\times_{S_j}S_i$, $G=G_j\times_{S_j}S$, and define $H_i$ and $H$ similarly. Then the canonical map below is bijective:
\[
\varinjlim_{i\geq j}\Hom_{S_i\text{-gr.}}(G_i,H_i)\to\Hom_{S\text{-gr.}}(G,H).
\]
\end{lemma}
\begin{lemma}{10.3}\label{VIB.10.3}
Let $G$ be an $S$-group of finite presentation; then there is $j\in I$, an $S_j$-group of finite presentation $G_j$, and an isomorphism of $S$-groups $G\cong G_j\times_{S_j}S$.
\end{lemma}
Assertions 10.2 and 10.3 are easy consequences of 10.0 and 10.1, taking account of the interpretation\nde{In terms of commutative diagrams of $S$-morphisms.} of the $S$-group structure given in EGA 0$_{\mathrm{III}}$, 8.2.5 and 8.2.6.

\begin{lemma}{10.4}\label{VIB.10.4}
Let $u:G\to H$ be a morphism of $S$-groups between $S$-groups of finite presentation. By 10.3 and 10.2, $u$ comes by base change from a morphism $u_j:G_j\to H_j$ between $S_j$-groups of finite presentation. Then, for $u$ to be a central monomorphism (resp. an invariant monomorphism), it is necessary and sufficient that there be $i\geq j$ such that $u_i=u_j\times_{S_j}S_i$ has the same property.
\end{lemma}
This is an immediate consequence of 10.0 and 10.1, taking account of the characterization given in 6.7 of central or invariant group monomorphisms.

\begin{corollary}{10.5}\label{VIB.10.5}
Let $G_j$ be an $S_j$-group of finite presentation. For $G_j\times_{S_j}S$ to be commutative, it is necessary and sufficient that there be $i\geq j$ such that $G_j\times_{S_j}S_i$ is so.
\end{corollary}
Indeed, it amounts to the same thing to say that an $S$-group is commutative, or that, considered as a subgroup scheme of itself, it is central.

\begin{proposition}{10.6}\label{VIB.10.6}
Let $G_j$ be an $S_j$-group of finite presentation, $G'_j$ a subgroup scheme of $G_j$ flat and of finite presentation over $S_j$. For $(G_j\times_{S_j}S)/(G'_j\times_{S_j}S)$ to be representable for the (fpqc) topology, it is necessary and sufficient that there be $i\geq j$ such that $(G_j\times_{S_j}S_i)/(G'_j\times_{S_j}S_i)$ is so.
\end{proposition}
This is a consequence of the following more general lemma:

\begin{lemma}{10.7}\label{VIB.10.7}
Let $X_j$ be an $S_j$-scheme of finite presentation, and $R_j$ an equivalence relation on $X_j$ flat and of finite presentation.\nde{I.e., such that the composite $R_j\hookrightarrow X_j\times_{S_j}X_j\xrightarrow{\mathrm{pr}_1}X_j$ is flat and of finite presentation.} For the quotient sheaf $(X_j\times_{S_j}S)/(R_j\times_{S_j}S)$ for the topology $T=\text{(fppf)}$ or $\text{(fpqc)}$ to be representable, it is necessary and sufficient that there be $i\geq j$ such that the quotient sheaf $(X_j\times_{S_j}S_i)/(R_j\times_{S_j}S_i)$ for the topology $T$ is so.
\end{lemma}
Taking account of the statements of EGA IV$_2$, 8.8.2, 8.8.3, 8.10.5 and 11.2.6 recalled in 10.0, this lemma is a consequence of the following result:
% typo?: kept as printed: EGA IV_2, rather than IV_3, in this reference.

\begin{lemma}{10.8}\label{VIB.10.8}
Let $T$ be the (fppf) or (fpqc) topology; let $X$ be an $S$-scheme of finite presentation (resp. locally of finite presentation), $R$ an equivalence relation on $X$ defined by a monomorphism $v:R\to X\times_S X$ such that $\mathrm{pr}_1\circ v$ is flat and of finite presentation (resp. flat and locally of finite presentation). Then the following conditions are equivalent:

(i) The quotient sheaf $X/R$ for the topology $T$ is representable.

(ii) There is an $S$-scheme of finite presentation (resp. locally of finite presentation) $Y$ and a faithfully flat morphism $p:X\to Y$ such that the diagram
\begin{equation*}\tag{D}
\xymatrix{R\ar[r]^{\mathrm{pr}_1\circ v}\ar[d]_{\mathrm{pr}_2\circ v}&X\ar[d]^p\\X\ar[r]_p&Y}
\end{equation*}
is cartesian.
\end{lemma}
Note first that by IV, 3.3.2 and 4.4.3, for the sheaf $X/R$ for the topology $T$ to be representable by $Y$, it is necessary and sufficient that the diagram (D) be cartesian and $p$ covering for the topology $T$.

Let us show that (i) implies (ii). Hypothesis (i) implies that the diagram (D) is cartesian, hence that $\mathrm{pr}_1\circ v$ is obtained from $p$ by base change by $p$, and that $p$ is covering for the (fpqc) topology. Thus, by (fpqc) descent (EGA IV$_2$, 2.7.1), since $\mathrm{pr}_1\circ v$ is faithfully flat and (locally) of finite presentation, the same holds for $p$. Then, by EGA IV$_4$, 17.7.5, since $X$ is (locally) of finite presentation over $S$, the same holds for $Y$.

Let us show that (ii) implies (i). It suffices to show that $p$ is covering for the (fppf) topology; now $p$ is faithfully flat by hypothesis, and is locally of finite presentation because $X$ and $Y$ are locally of finite presentation over $S$ (EGA IV$_1$, 1.4.3 (v)).

\begin{lemma}{10.9}\label{VIB.10.9}
Let $G_j$ be an $S_j$-group of finite presentation, and $G=G_j\times_{S_j}S$. For $G^0$ to be representable, it is necessary and sufficient that there be $i\geq j$ such that $(G_i)^0=(G_j\times_{S_j}S_i)^0$ is so.
\end{lemma}
The condition is sufficient, since the functor $G\mto G^0$ commutes with base change, by 3.3.

Conversely, suppose $G^0$ representable. Then, by 3.9, $G^0$ is open in $G$ and quasi-compact over $S$, hence of finite presentation over $S$, since $G$ is. Then, by 10.3 and 10.1, there is $i\geq j$ and an open subgroup scheme $H_i$ of $G_i$ such that $H_i\times_{S_i}S=G^0$. The structural morphism $G^0\to S$ is connected, i.e.\ has geometrically connected fibers (VIA 2.1.1), so, by EGA IV$_3$, 9.3.3 and 9.7.7, increasing $i$ if necessary, one may suppose that the structural morphism $H_i\to S$ is connected. Then, by 3.10.1, the underlying space of $H_i$ is precisely $\underline{(G_i)^0}$, and thus $H_i$ represents $(G_i)^0$.
% typo?: kept as printed: H_i -> S, rather than H_i -> S_i.

\numpara{10.10}\label{VIB.10.10}
Recall two very useful special cases of the situation stated in 10.0 (cf. EGA IV$_3$, 8.1.2 a) and c)):

a) Given a point $x$ of a scheme $X$, one puts $S_0=\Spec\ZZ$ and considers the decreasing filtered projective system $(S_i)_{i\in I}$ of affine open neighborhoods of $x$; then $S=\Spec\OO_{X,x}$. In particular, if $x$ is the generic point of an integral scheme $X$, one finds $S=\Spec\kappa(x)$.

b) One puts $S_0=\Spec\ZZ$, and considers the family $(\mathscr A_i)_{i\in I}$, preordered by inclusion, of sub-$\ZZ$-algebras of finite type of the ring of an affine scheme $S$. Since the $\mathscr A_i$ are noetherian rings, this allows one in many cases to pass from the noetherian case to the general case.

We shall now give two results concerning the special case considered in a).\nde{Note that the proof also uses case b).}

\begin{proposition}{10.11}\label{VIB.10.11}
\nde{The statement has been simplified, and the case of subgroups treated separately in Corollary 10.11.1.}
Let $S$ be an integral scheme with generic point $\eta$, $G$ (resp. $Y$, $Z$) an $S$-group (resp. $S$-schemes) of finite presentation, $u,v,i:Y\to G$ and $j:Z\to G$ morphisms of $S$-schemes. Suppose that $i_\eta:Y_\eta\to G_\eta$ and $j_\eta:Z_\eta\to G_\eta$ are closed immersions.

Then there is a nonempty open subset $S'$ of $S$ such that the morphisms $i':Y'\to G'$ and $j':Z'\to G'$ obtained by base change are closed immersions, and the functors:
\[
\uTransp(u',v'),\quad\uTransp_{G'}(i'(Y'),j'(Z'))\quad\text{and}\quad\uTransp^{\mathrm{str}}_{G'}(i'(Y'),j'(Z'))
\]
resp.
\[
\uCentr(u')\quad\text{and}\quad\uNorm_{G'}i'(Y')
\]
are representable by closed sub-$S'$-schemes (resp. sub-$S'$-groups) of $G'$, of finite presentation over $S'$.
\end{proposition}
We shall apply the results of 10.1, first in the situation of 10.10 a), then in that of 10.10 b). Since $G_\eta$, $Y_\eta$, $Z_\eta$ are flat over the field $\kappa(\eta)$, $G_\eta$ is separated over $\kappa(\eta)$ (VIA 0.3), and $i_\eta$, $j_\eta$ are closed immersions, then, by 10.1, there is an affine open subset $S'=\Spec(A')$ of $S$, a noetherian subring $A$ of $A'$, $A$-schemes $G_A$, $Y_A$, $Z_A$ flat and of finite presentation over $A$, and morphisms $u_A,v_A,i_A:Y_A\to G_A$ and $j_A:Z_A\to G_A$, such that $G_A$ is an $A$-group, separated over $A$, $i_A$ and $j_A$ are closed immersions, and $G\times_S S'=G_A\otimes A'$, etc. Since the functors considered for $S'$ are obtained by base change from the analogous functors for $\Spec(A)$, it suffices to establish the result for the latter.

By EGA IV$_2$, 6.9.2, replacing $A$ by a localization $A_f$ if necessary (and hence $S'$ by the affine open subset $S'_f$), one may suppose that $G_A$, $Y_A$, $Z_A$ are essentially free over $A$ (in the sense of 6.2.1).\nde{Indeed, since $G_A$ is flat and of finite presentation over $A$, it is covered by affine open subsets $G_1,\ldots,G_n$ such that each $\OO(G_i)$ is a flat $A$-algebra of finite presentation; then, by EGA IV$_2$, 6.9.2, there is $f_i\in A$ such that $\OO(G_i)_{f_i}$ is a free module over $A_{f_i}$; one can then replace $\Spec(A)$ by the affine open subset $D(f)$, where $f=f_1\cdots f_n$, and one does the same for $Y_A$ and $Z_A$.} It then follows from 6.2.4 b) and e) that, under the hypotheses of the statement, the functors considered are representable by closed sub-$A$-schemes of $G_A$ (hence of finite presentation over $A$, since $A$ is noetherian and $G_A$ of finite presentation over $A$), and they are sub-$A$-groups of $G_A$ in the case of $\uCentr(u_A)$ and $\uNorm_{G_A}i_A(Y_A)$.
